\section{Обзор литературы и теоретические основы}

\subsection{Гидродинамическая теория смазки}

Основой анализа подшипников скольжения является гидродинамическая теория смазки, заложенная в работах О.~Рейнольдса~\cite{pinkus}. Несущая способность смазочного слоя возникает вследствие вязкого течения жидкости в клиновом зазоре между вращающимся валом и неподвижной втулкой.

Основные допущения~\cite{hamrock}: смазка~-- несжимаемая ньютоновская жидкость; течение ламинарное; инерционные силы пренебрежимо малы; толщина смазочного слоя мала по сравнению с радиусом ($c \ll R$); давление постоянно по толщине плёнки; отсутствует проскальзывание на границах.

\subsection{Дискретный микрорельеф поверхностей трения}

Технология лазерного текстурирования поверхностей получила широкое развитие~\cite{etsion2005}. Основные механизмы влияния микрорельефа~\cite{gropper2016, fowell2007}:

\textbf{Микрогидродинамический эффект.} Каждое углубление создаёт локальный клиновый зазор с дополнительным гидродинамическим давлением.

\textbf{Снижение сдвиговых напряжений.} Увеличение средней толщины смазочного слоя уменьшает градиент скорости и коэффициент трения.

\textbf{Резервуары смазки.} Углубления обеспечивают подачу смазки при пуске и останове~\cite{galda2009}.

Эффективность текстурирования зависит от формы, глубины, размеров и плотности расположения углублений~\cite{tala2012, brizmer2003, rahmani2007}.

\newpage
